\section{Comunicación y animación con Python}

Una vez modelados la planta física y los tres niveles de control, se integran todos los
subsistemas en un único modelo de Simulink que constituye el sistema completo a simular. En él se
interconectan la planta con las capas de control regulatorio, supervisor y de
protección, junto con el bloque de comunicación que enlaza la simulación con la
animación. La Fig. \ref{fig:planta_completa} muestra esta integración general.

\begin{figure}[H]
    \centering
    \includegraphics[width=0.85\textwidth]{simulacion/planta_completa.png}
    \caption{Modelo completo del sistema en Simulink: planta física e integración de los tres niveles de control.}
    \label{fig:planta_completa}
\end{figure}

Para la visualización en tiempo real de la grúa y su entorno, se implementó un bloque de comunicación UDP que transmite los estados
críticos del sistema a un programa de python que lo comunica con la simulación. Este programa muestra en una ventana del navegador 
una representación 2D de la grúa, contenedores, barco y parámetros clave que permiten al operador monitorear visualmente la operación.

\begin{figure}[H]
    \centering
    \includegraphics[width=0.3\textwidth]{simulacion/sim_udp.png}
    \hfill
    \includegraphics[width=0.3\textwidth]{simulacion/sim_udp2.png}
    \caption{Diagrama de la comunicación UDP entre Simulink y el programa de python}
\end{figure}

A continuación se presentan algunas capturas de diferentes estados de la grúa, mostrando la posición del spreader, la carga y el entorno.

\begin{figure}[H]
    \centering
    \begin{subfigure}{0.49\textwidth}
        \centering
        \includegraphics[width=\linewidth]{simulacion/sim_manual.png}
        \caption{Modo manual.}
    \end{subfigure}
    \hfill
    \begin{subfigure}{0.49\textwidth}
        \centering
        \includegraphics[width=\linewidth]{simulacion/sim_auto.png}
        \caption{Modo automático.}
    \end{subfigure}

    \vspace{0.4cm}

    \begin{subfigure}{0.49\textwidth}
        \centering
        \includegraphics[width=\linewidth]{simulacion/sim_auto_leave.png}
        \caption{Modo automático: liberación de la carga.}
    \end{subfigure}
    \hfill
    \begin{subfigure}{0.49\textwidth}
        \centering
        \includegraphics[width=\linewidth]{simulacion/sim_error.png}
        \caption{Estado de error.}
    \end{subfigure}

    \caption{Capturas de la animación en distintos estados de operación de la grúa.}
    \label{fig:sim_estados}
\end{figure}

